
@@M0@@) @@M1@@
sont linéairement dépendants sur @@M2@@.

@@M3@@) @@M4@@ sont linéairement indépendants sur @@M5@@.

1.2.
Pour plus de commodité, nous allons rappeler ici quelques définitions et propriétés classiques des Modules gradués.

Soient @@M6@@ un schéma, @@M7@@ une @@M8@@-Algèbre graduée, @@M9@@ un @@M10@@-Module gradué, Pour tout @@M11@@-Module gradué @@M12@@, on désigne comme d'habitude par @@M13@@ le @@M14@@-Module gradué défini par :
@@M15@@
On dit que @@M16@@ est un @@M17@@-Module gradué localement libre (resp. localement libre de type fini) si tout point de @@M18@@ possède un voisinage ouvert @@M19@@ tel que @@M20@@ soit isomorphe, en tant que @@M21@@-Module gradué, à une somme directe de Modules de la forme @@M22@@ (resp. une somme directe finie).

On dit que @@M23@@ est un @@M24@@-Module gradué de présentation finie si tout point de @@M25@@ possède un voisinage ouvert @@M26@@ tel qu'il existe une suite exacte @@M27@@, où @@M28@@ et @@M29@@ sont des @@M30@@-Modules gradués libres de type fini, et les homomorphismes compatibles aux graduations.

1.2.1. Pour tout entier @@M31@@, on peut définir un homomorphisme canonique de @@M32@@-Modules gradués :
@@M33@@
en envoyant @@M34@@, où @@M35@@ est une section de @@M36@@, et @@M37@@ une section de @@M38@@ sur l'élément @@M39@@ de @@M40@@.

1.2.2. On dit qu'un @@M41@@-Module gradué @@M42@@ est libre de type fini, et homogène de degré @@M43@@ s'il est isomorphe en tant que @@M44@@-Module gradué à @@M45@@ pour un certain entier @@M46@@.

